\documentclass[14pt]{beamer}
\usepackage{pepslides}
\title{Self-gravitation and rotation}
\lecturenumber{21}
\hypersetup{pdftitle={Lecture 21: Self-gravitation and rotation}}

\begin{document}

\titleslide

\begin{frame}[t]
\heading{From earlier lectures}
\centering
\begin{minipage}[b]{0.5\textwidth}
\centering
\includegraphics[width=\textwidth]{L6F5.pdf}\\[0.5em]
Preliminary Reference Earth Model (PREM)\\ of Dziewonski \& Anderson (1981)
\end{minipage}\hfill
\begin{minipage}[b]{0.42\textwidth}
\centering
\includegraphics[height=0.6\paperheight]{L8F6.pdf}\\[0.5em]
GLAD-M35 of Cui \emph{et al.} (2024) -- \\ a global waveform inverse model
\end{minipage}
\end{frame}

\sectionslide{Newtonian gravity}

\begin{frame}
Gravitational potential satisfies \textbf{Poisson's equation}:
\begin{eqn}
\nabla^{2}\phi=4\pi G\varrho.
\end{eqn}
Associated gravitational acceleration:
\begin{eqn}
g_{i}=-\frac{\partial\phi}{\partial x_{i}}.
\end{eqn}
Equation subject to continuity conditions:
\begin{eqn}
[\phi]_{-}^{+}=0,\quad [g_{i}\hat{n}_{i}]_{-}^{+}=0,
\end{eqn}
across any jumps in density, while the potential also tends to zero at infinity.
\end{frame}

\begin{frame}
Solution to the problem can be expressed as a familiar integral:
\begin{eqn}
\phi(\mathbf{y},t)=-G\int_{M_{t}}\frac{\varrho(\mathbf{y}',t)}{[(y_{i}-y_{i}')(y_{i}-y_{i}')]^{\frac{1}{2}}}\dd^{3}\mathbf{y}',
\end{eqn}
and differentiating, the gravitational acceleration can be similarly written:
\begin{eqn}
g_{i}(\mathbf{y},t)=-G\int_{M_{t}}\frac{\varrho(\mathbf{y}',t)(y_{i}-y_{i}')}{[(y_{j}-y_{j}')(y_{j}-y_{j}')]^{\frac{3}{2}}}\dd^{3}\mathbf{y}'.
\end{eqn}
\end{frame}

\begin{frame}
\centering
\includegraphics[height=0.62\paperheight]{L9F1.pdf}

\raggedright
Gravitational acceleration and pressure in PREM, calculated from its density profile.

The gravitational acceleration is close to $10\,\mathrm{m\,s^{-2}}$ throughout the mantle, and largest at the core--mantle boundary.
\end{frame}

\begin{frame}
Using the motion of the body, we can define a \textbf{referential gravitational potential}:
\begin{eqn}
\zeta(\mathbf{x},t)=\phi[\bphi(\mathbf{x},t),t].
\end{eqn}
Using conservation of mass, we can transform the integral solution to obtain:
\begin{eqn}
\zeta(\mathbf{x},t)=-G\int_{M}\frac{\rho(\mathbf{x}')}{\{[\varphi_{i}(\mathbf{x},t)-\varphi_{i}(\mathbf{x}',t)]
[\varphi_{i}(\mathbf{x},t)-\varphi_{i}(\mathbf{x}',t)]\}^{\frac{1}{2}}}\dd^{3}\mathbf{x}'.
\end{eqn}
Similarly for the \textbf{referential gravitational acceleration}:
\begin{eqn}
\gamma_{i}(\mathbf{x},t)=-G\int_{M}\frac{\rho(\mathbf{x}')[\varphi_{i}(\mathbf{x},t)-\varphi_{i}(\mathbf{x}',t)]}
{\{[\varphi_{j}(\mathbf{x},t)-\varphi_{j}(\mathbf{x}',t)][\varphi_{j}(\mathbf{x},t)-\varphi_{j}(\mathbf{x}',t)]\}^{\frac{3}{2}}}\dd^{3}\mathbf{x}'.
\end{eqn}
\end{frame}

\begin{frame}
The \textbf{gravitational binding energy} is the potential energy of a body associated with gravitational forces.

It can be calculated by assuming the body is assembled bit-by-bit, summing up the work
required to move each particle in from infinity relative to the current field.

Calculation identical to that in Part IB Electrodynamics for the energy of an electrostatic field:
\begin{eqn}
\mathcal{V}_{g}(t)=\frac{1}{2}\int_{M_{t}}\varrho(\mathbf{y},t)\phi(\mathbf{y},t)\dd^{3}\mathbf{y}.
\end{eqn}
Full details in the notes, but you won't be asked to recall the calculation.

Can express referentially:
\begin{eqn}
\mathcal{V}_{g}(t)=\frac{1}{2}\int_{M}\rho(\mathbf{x})\zeta(\mathbf{x},t)\dd^{3}\mathbf{x}.
\end{eqn}
\end{frame}

\sectionslide{Equations of motion}

\begin{frame}
Potential energy is now the sum of elastic and gravitational contributions.

Modified Lagrangian density:
\begin{eqn}
\mathcal{L}(\mathbf{x},t,\bphi,\mathbf{v},\mathbf{F})=\frac{1}{2}\rho(\mathbf{x})\|\mathbf{v}\|^{2}
-W(\mathbf{x},t,\mathbf{F})-\frac{1}{2}\rho(\mathbf{x})\zeta(\mathbf{x},t).
\end{eqn}
Statement of Hamilton's principle is unchanged.

Derivation of the Euler--Lagrange equations is harder because we have a \textbf{non-local} term in
the Lagrangian -- this reflects the \textbf{action-at-a-distance} property of Newtonian gravity.
\end{frame}

\begin{frame}
Functional derivatives are still defined through the linear term in a Taylor expansion.

The result can be expressed in the form
\begin{eqn}
\langle D\mathcal{S}(\bphi)|\delta\bphi\rangle=\int_{0}^{T}\int_{M}
\left(\frac{\delta\mathcal{L}}{\delta\varphi_{i}}\delta\varphi_{i}+\frac{\delta\mathcal{L}}{\delta v_{i}}
\delta v_{i}+\frac{\delta\mathcal{L}}{\delta F_{ij}}\delta F_{ij}\right)\dd^{3}\mathbf{x}\dd t.
\end{eqn}
The ``derivative'' terms here are, by definition, just the quantities multiplying the first-order
perturbations to the fields.

In the case of local Lagrangians, these are just normal partial derivatives.

It is only the variation with respect to the motion that needs additional work.
\end{frame}

\begin{frame}
Following the standard integrations by parts, we obtain:
\begin{eqn}
\begin{gathered}
\frac{\delta\mathcal{L}}{\delta\varphi_{i}}-\frac{\partial}{\partial t}\left(\frac{\delta\mathcal{L}}
{\delta v_{i}}\right)-\frac{\partial}{\partial x_{j}}\left(\frac{\delta\mathcal{L}}{\delta F_{ij}}\right)=0,\\[\eqgap]
\frac{\delta\mathcal{L}}{\delta F_{ij}}\hat{n}_{j}=0.
\end{gathered}
\end{eqn}
As previously, we have:
\begin{eqn}
\frac{\delta\mathcal{L}}{\delta v_{i}}=\rho v_{i},\quad \frac{\delta\mathcal{L}}{\delta F_{ij}}=-\frac{\partial W}{\partial F_{ij}}.
\end{eqn}
In the second problem set you will show that:
\begin{eqn}
\frac{\delta\mathcal{L}}{\delta\varphi_{i}}=\rho\gamma_{i}.
\end{eqn}
\end{frame}

\begin{frame}
The equations of motion can be written in the form:
\begin{eqn}
\rho\frac{\partial v_{i}}{\partial t}-\frac{\partial T_{ij}}{\partial x_{j}}-\rho\gamma_{i}=0.
\end{eqn}
As expected, we have just added a gravitational body force to the balance of linear
momentum.

This force is expressed in terms of the motion through an integral over the whole body.

These are now \textbf{integro-differential equations}.
\end{frame}

\begin{frame}
\begin{columns}[c]
\begin{column}{0.4\textwidth}
\centering
\includegraphics[width=\textwidth]{dehant2015_rotation.jpg}
\end{column}
\begin{column}{0.56\textwidth}
We have described the motion relative to an inertial frame defined by the ``fixed stars''.

For applications to the Earth, it is more natural to use a rotating frame.

The Earth's rotation is not steady, varying due to both external and internal processes.

Within seismology, it is common to use a steady and approximately co-rotating reference frame.

This will introduce the usual fictitious forces into the equations of motion.
\end{column}
\end{columns}
\credit{Dehant \& Mathews (2015)}
\end{frame}

\begin{frame}
The motion relative to inertial space can be written:
\begin{eqn}
\varphi_{i}(\mathbf{x},t)=\Phi_{i}(t)+R_{ij}(t)\overline{\varphi}_{j}(\mathbf{x},t).
\end{eqn}
The translation vector represents motion of the origin of the reference frame.

The rotation matrix represents rotation of the reference frame.

On the right-hand side we see the \textbf{relative motion} as seen from the non-inertial frame.

Because of material frame indifference, the choice of frame has no effect on the potential
energy associated with the motion.
\end{frame}

\begin{frame}
Differentiating the motion in time, we find:
\begin{eqn}
v_{i}=\frac{\ddns \Phi_{i}}{\ddns t}+R_{ij}\left(\frac{\partial\overline{\varphi}_{j}}{\partial t}+R_{kj}\frac{\ddns R_{kl}}{\ddns t}\overline{\varphi}_{l}\right).
\end{eqn}
Introduction of the angular velocity as seen from the non-inertial frame:
\begin{eqn}
R_{kj}\frac{\ddns R_{kl}}{\ddns t}=\epsilon_{jml}\Omega_{m}.
\end{eqn}
The velocity can then be expressed in the form:
\begin{eqn}
v_{i}=\frac{\ddns \Phi_{i}}{\ddns t}+R_{ij}(\overline{v}_{j}+\epsilon_{jkl}\Omega_{k}\overline{\varphi}_{l}).
\end{eqn}
\end{frame}

\begin{frame}
Differentiating again, we can write the acceleration:
\begin{eqn}
\begin{aligned}
\frac{\partial v_{i}}{\partial t}=\frac{\ddns ^{2}\Phi_{i}}{\ddns t^{2}}+R_{ij}\biggl(&\frac{\partial\overline{v}_{j}}{\partial t}+2\epsilon_{jkl}\Omega_{k}\overline{v}_{l}\\
&+\epsilon_{jkl}\frac{\ddns \Omega_{k}}{\ddns t}\overline{\varphi}_{l}+\epsilon_{jkl}\epsilon_{lmn}\Omega_{k}\Omega_{m}\overline{\varphi}_{n}\biggr).
\end{aligned}
\end{eqn}
At this point, we define an approximately co-rotating frame by setting:
\begin{eqn}
\frac{\ddns \bm{\Phi}}{\ddns t}=\mathbf{0},\quad \frac{\ddns \bm{\Omega}}{\ddns t}=\mathbf{0}.
\end{eqn}
This approach can be generalised to allow for interactions between different planets
(e.g.\ the Earth--Moon system).
\end{frame}

\begin{frame}
Differentiating the motion spatially, we find:
\begin{eqn}
F_{ij}=R_{ik}\overline{F}_{kj}.
\end{eqn}
From the definition of the first Piola--Kirchhoff stress and frame indifference, we find:
\begin{eqn}
T_{ij}=R_{ik}\overline{T}_{kj}.
\end{eqn}
Similarly, for the referential gravitational acceleration:
\begin{eqn}
\gamma_{i}=R_{ij}\overline{\gamma}_{j}.
\end{eqn}
\end{frame}

\begin{frame}
Putting things together and cancelling common rotation factors:
\begin{eqn}
\rho\left(\frac{\partial\overline{v}_{i}}{\partial t}+2\epsilon_{ijk}\Omega_{j}\overline{v}_{k}+\epsilon_{ijk}
\epsilon_{klm}\Omega_{j}\Omega_{l}\overline{\varphi}_{m}\right)-\frac{\partial\overline{T}_{ij}}{\partial x_{j}}-\rho\overline{\gamma}_{i}=0.
\end{eqn}
All terms are defined relative to the rotating reference frame.

Identical to the inertial case except for the addition of \textbf{Coriolis and centrifugal forces}.

\textbf{From here on we will always work in the rotating frame and drop the overbars.}
\end{frame}

\begin{frame}
Same procedure as ever for linearisation on application of a small stress glut.

At zeroth order we have the equilibrium equations:
\begin{eqn}
\rho\epsilon_{ijk}\epsilon_{klm}\Omega_{j}\Omega_{l}x_{m}-\frac{\partial T_{ij}^{0}}{\partial x_{j}}-\rho\gamma_{i}^{0}=0.
\end{eqn}
\textbf{Stress-free equilibrium solutions are not now possible.}
\end{frame}

\begin{frame}
At first order we have the linearised equations for the displacement vector:
\begin{eqn}
\begin{aligned}
&\rho\left(\frac{\partial^{2}u_{i}}{\partial t^{2}}+2\epsilon_{ijk}\Omega_{j}\frac{\partial u_{k}}{\partial t}+
\epsilon_{ijk}\epsilon_{klm}\Omega_{j}\Omega_{l}u_{m}\right)\\
&\qquad-\frac{\partial}{\partial x_{j}}\left(A_{ijkl}
\frac{\partial u_{k}}{\partial x_{l}}\right)-\rho\gamma_{i}^{1}=-\frac{\partial\overline{S}_{ij}}{\partial x_{j}}.
\end{aligned}
\end{eqn}
The perturbation to the referential gravitational acceleration is given as an integral of the
displacement vector -- see the handout for the full expression.
\end{frame}

\sectionslide{Scaling analysis}

\begin{frame}
Derivatives can be readily estimated in terms of typical time and length scales:
\begin{eqn}
\begin{aligned}
\rho\frac{\partial^{2}u_{i}}{\partial t^{2}}&\sim\frac{\overline{\rho}U}{T^{2}},\\[0.4em]
\rho\epsilon_{ijk}\Omega_{j}\frac{\partial u_{k}}{\partial t}&\sim\frac{\overline{\rho}\Omega U}{T},\\[0.4em]
\rho\epsilon_{ijk}\epsilon_{klm}\Omega_{j}\Omega_{l}u_{m}&\sim\overline{\rho}\Omega^{2}U,\\[0.4em]
\frac{\partial}{\partial x_{j}}\left(A_{ijkl}\frac{\partial u_{k}}{\partial x_{l}}\right)&\sim\frac{\overline{\mu}U}{L^{2}}.
\end{aligned}
\end{eqn}
Estimating the magnitude of the gravitational acceleration is a bit more work due to its
integral definition (see handout):
\begin{eqn}
\rho\gamma_{i}^{1}\sim4\pi G\overline{\rho}^{2}U.
\end{eqn}
\end{frame}

\begin{frame}
The ratios of the terms, with $L$ a wavelength and $T$ a period:
\begin{eqn}
\frac{\textrm{gravitational forces}}{\textrm{elastic forces}}=\frac{4\pi G\overline{\rho}L^{2}}{\overline{v}^{2}},\qquad
\frac{\textrm{Coriolis term}}{\textrm{inertial term}}=\Omega T.
\end{eqn}
\begin{center}
\begin{tabular}{llcc}
\toprule
period & & gravity/elasticity & Coriolis/inertia\\
\midrule
1 s & short-period body waves & $5\times10^{-6}$ & $7\times10^{-5}$\\
20 s & long-period body waves & $2\times10^{-3}$ & $1\times10^{-3}$\\
100 s & surface waves & $0.05$ & $0.007$\\
1 hour & longest free oscillations & $3$ & $0.3$\\
12 hours & tides & $3$ & $3$\\
\bottomrule
\end{tabular}
\end{center}
{\small With $L=\overline{v}T$, but no larger than the radius of the Earth, $\overline{v}\approx8\,\mathrm{km\,s^{-1}}$ and
$\overline{\rho}\approx5.5\,\mathrm{g\,cm^{-3}}$.\par}

\textbf{At the longest wavelengths gravity is just as important as elasticity within the
equations of motion -- this fact means that we can learn about density directly.}
\end{frame}

\begin{frame}
\centering
\includegraphics[height=0.74\paperheight]{L9F2.pdf}

\raggedright
Synthetic seismograms in PREM with gravitation (black) and with all gravitational effects switched off (red),
in successively lower frequency bands.
\end{frame}

\begin{frame}[t]
\heading{What you need to know and be able to do}
\begin{enumerate}
\item[(i)] Understand the relevant parts of Newtonian gravity, including the manner in which the gravitational
potential and acceleration can be transformed into quantities defined on the reference body.
\item[(ii)] Know physically what the gravitational binding energy of a planet represents, and how it can be
determined.
\item[(iii)] Know how the gravitational binding energy is incorporated into the Lagrangian, and how the equations of
motion can be obtained. The full derivation is too long to be examined, but you may be asked to reproduce parts of it
given appropriate information and guidance.
\item[(iv)] Perform scaling arguments to determine the relative size of terms within an equation, and to determine
whether some of them might be negligible. This is a technique required throughout this course, and you should ensure
you are happy with it.
\end{enumerate}
\end{frame}

\end{document}
